\silentchapter{Appell}{appell}
\begin{widebody}
{\small

Cellekolonien som legg seg til ro på bunnen av petriskåla og avgjer om ho skal verte eit kollektiv eller ikkje, utfører same operasjon som eit frogeembryo som legg ansiktet sitt saman frå eit bioelektrisk minne ingen einskild celle rår over. Det elektriske ansiktet; kartlagt i fluorescerande fargestoff av Levins laboratorium\textsuperscript{22}, synleg som eit mønster av spenningsgradientar lenge før gena som kodar for auge og munn er slått på; er eit minne om noko som ikkje finst enno, bore av tusenvis av celler som har semjast om ein felles modell av det dei skal verte. Éin organisme finst fordi alle cellene er samde om at det er dit dei skal. Skrap ein strek i plastret, og kvar øy bestemmer seg for at ho er embryoet. Slik oppstår tvillingar. Talet på individ i eit blastoderm er ikkje genetisk bestemt; det er eit fysiologisk utfall av kor langt den bioelektriske samordninga rekk.

Kreft, i dette rammeverket, er ikkje ein sjukdom i cellene. Det er ein sjukdom i sjølvet. Onkogenet koplar cella elektrisk frå nettverket. Grensa mellom sjølv og omverd krympar. Cella er ikkje meir egoistisk; ho har eit mindre sjølv. Ho sluttar å delta i det kollektive prosjektet og vert ein eincelleorganisme omgjeven av eit miljø ho ikkje lenger vedkjenner seg. Metastase. Men tvingar du cella tilbake inn i det elektriske nettverket, forsvinn svulsten. Onkoproteinet glør framleis raudt i kvart fluorescensmikroskopbilete. Genfeilen er der. Men cella navigerer att mot anatomiske mål ho hadde gløymt fordi ho var kopla frå nettverket som bar dei.

Eg skriv dette til deg som formar.

Det du nettopp las er ikkje ein metafor for design. Det er same struktur. Kvar gong du snevrar \textbf{lyskjegla}\textsuperscript{21}; kvar gong du reduserer alle \textbf{seleksjonstrykk} til eitt; kvar gong du optimaliserer éin dimensjon og let resten fare; gjer du det same med forma som onkogenet gjer med cella. Du koplar henne frå nettverket av agentar som saman heldt forma stabil. Materialet mistar stemma si. Handverkaren mistar hendene sine. Brukaren mistar kroppen sin. Økosystemet mistar framtida si. Forma vert ein amøbe; ho sviktar langs aksane ingen representerte, og svikten vert boren av dei som lever med henne, i åra og tiåra etter at du har levert teikningane og fått honoraret og gått vidare til neste prosjekt.

Trachealceller frå eit menneske, frigjorde frå kroppen, reorganiserer seg til sjølvdrivne vesnar som aldri har eksistert på jorda, og byrjar lækje sår ingen har bede dei lækje. Halvparten av genomet vert uttrykt annleis. Ni tusen gen skiftar tilstand. Same DNA, radikalt annleis form. Genregulatoriske nettverk, berre kjemiske stoff i samspel, er i stand til seks ulike former for læring. Ein sorteringsalgoritme, deterministisk og triviell, har sideoppdrag algoritmen aldri bad om; aspekt av kognisjon, minimale men reelle, heilt ned til botn. Det finst ikkje passiv materie. Det har aldri gjort det. Det me kalla dødt stoff var levande kompetanse me ikkje hadde lyskjegle til å sjå.

Og me byggjer no. Me byggjer med agentisk materie. Me byggjer med språkmodellar som navigerer latente rom. Me byggjer med genredigering og bioprinting og sjølvmonterande material og fleirskala-kompetansehierarki der celle, materiale, verktøy, handverkar, organisasjon og marknad alle er agentar med eigne lyskjegler som deformerer landskapet for kvarandre samstundes. Me byggjer med teknologiar som gjer det mogleg å koordinere kompetanse på skalaer som aldri har vore kopla; frå molekylær manipulasjon til kontinental økologisk restaurering.

Me byggjer utan eit fundament.

Designfaga har ingen formell modell for kva dei gjer. Ingen felles grammatikk. Ingen presis definisjon av kva ein agent er, kva eit \textbf{seleksjonstrykk} er, kva eit tilpassingslandskap er, kva omsorga til ein formgjevar faktisk rommar. Me har estetisk intuisjon. Me har stilhistorie. Me har handverksminne. Det er ikkje nok. Det har aldri vore nok. Men konsekvensane av å mangle det har aldri vore so store, fordi verktøya aldri har vore so mektige, og materien aldri so tydeleg i stand til å svare tilbake.

Kvart materialeøkosystem me kollapsar, kvar handverkskompetanse me fortrenger, kvart lokalt agenthierarki me faldar til éin global attraktor fordi éin dimensjon er lettare å rekne på enn fem; det er celler som vert kopla frå nettverket. Kompetansane let seg ikkje gjenskape identisk frå læra som fortrengte dei. Tapet er irreversibelt. Me krympar sjølvet til formverda kvar einaste gong me snevrar \textbf{lyskjegla}\textsuperscript{21} til det som er enkelt å optimalisere.

Denne traktaten er ikkje ei sanning. Ho er eit koordinatsystem. Ho tvingar deg til å seie kva du representerer og kva du overlet. Ho tvingar deg til å spesifisere kva agentane er, kva \textbf{seleksjonstrykk}a er, korleis landskapet ser ut, kor vid lyskjegla di er. Og det du ikkje kan spesifisere, kan du ikkje forsvare. Rammeverket er falsifiserbart; vilkåra for å felle det er lista i notata. Om du kan felle det, fell det. Det ville vere eit framsteg.

Men inntil det fell: bruk det. Utvid lyskjegla. Representér fleire aksar. Kopla deg tilbake til nettverket av agentar du formar saman med, ikkje over. Bygg former som toler kreftene du ikkje kontrollerer. Ikkje fordi traktaten seier det. Men fordi designhandlinga har konsekvensar for andre agentar, og konsekvensane er borna av dei som lever med forma.

Aksolotlet som misser eit lem detekterer avviket frå måltilstanden, arbeider seg tilbake, og stansar når lemmet er korrekt. Stansinga er det mest forbløffande. Systemet veit kva det skal vere. Spørsmålet er om du veit kva formene dine skal vere; ikkje kva dei skal sjå ut som, men kor mange aksar dei skal tole, kor mange agentar dei skal tene, kor mange krefter dei skal stå i.
}
\end{widebody}
